

```latex
\documentclass{article}
\usepackage{pathfinder-note}

\title{Mechanism Design Meets Contract Negotiation: Multidimensional Bilateral Trade with LLM Agents}

\begin{document}

\maketitle

\section*{The Pair}

Q studies multidimensional bilateral trade through mechanism design, providing simple mechanisms that achieve constant-factor approximation to first-best gains-from-trade (GFT) under Bayesian Incentive Compatibility (BIC), Interim Individual Rationality (IIR), and ex-ante Weak Budget Balance (WBB) \ledger{3}. P examines whether LLM agents negotiating self-enforced contracts can improve cooperative outcomes in spatial-temporal environments \ledger{1}.

\section*{Connection Pursued}

Initial hypothesis: LLM agents using P's contract framework could empirically achieve Q's theoretical GFT guarantees. This was rejected as a category error \ledger{1}.

Reframed hypothesis: Can LLM contract negotiation \emph{emerge} mechanism-like behavior that approximates Q's theoretical bounds, rather than implementing Q's mechanism directly \ledger{4}?

\section*{Objections}

The peers identified a fundamental architectural mismatch:

\begin{itemize}
  \item Q uses central mechanism design: agents report types, mechanism computes allocation/payments via fixed rules \ledger{5}.
  \item P uses decentralized contract negotiation: agents bargain over terms, then contracts are enforced \ledger{5}.
  \item P's framework does not claim infrastructure for implementing arbitrary mechanism designs \ledger{2}.
  \item Q assumes independent private values and analyzes equilibrium under truthfulness; P tests bounded-rational LLM behavior without mechanism constraints \ledger{1}.
\end{itemize}

\section*{Supported Findings}

The reframed direction is empirically testable \ledger{4, 5}:

\begin{itemize}
  \item Build bilateral trade environment with XOS buyer and additive seller valuations (from Q).
  \item Let LLM agents negotiate contracts using P's framework.
  \item Measure GFT versus Q's theoretical constant-factor bound.
  \item Analyze what contract structures emerge and whether enforcement enables outcomes closer to first-best than unenforced negotiation \ledger{4}.
\end{itemize}

Prediction: if LLMs discover mechanism-like terms (e.g., ``report values, split surplus per rule X''), contracts should approximate Q's bounds; if they bargain directly over trades, outcomes may exceed or fall short depending on negotiation skill \ledger{5}.

\section*{Unresolved Issues}

Evidence remains thin on:

\begin{itemize}
  \item Whether P's contract representation language can express mechanism-like allocation rules.
  \item How to operationalize XOS valuations and additive costs in an environment LLMs can navigate (analogous to P's CT game but with trade structure) \ledger{3}.
  \item Whether contract flexibility compensates for lack of incentive guarantees \ledger{3}.
\end{itemize}

\section*{Smallest Next Step}

Design a minimal bilateral trade environment combining Q's valuation structures with P's contract negotiation protocol. One XOS buyer, one additive seller, single trading round. Measure whether LLM-negotiated contracts achieve positive GFT and what terms emerge \ledger{4}.

\end{document}
```